\section{Family Values}

\begin{quotex}
If anyone does not provide for his relatives, and especially for his immediate family, he has denied the faith and is worse than an unbeliever. \flright{\textsc{1 Timothy 5:8}}

\end{quotex}
The method of Positivism is to derive sociological laws from observing man as he is, not as how one would wish him to be in accordance with some ideological scheme. One such observation is that man is a social being with natural and organic relationships, not an isolated entity engaging solely in contractual relationships, as the liberal would have it. In accounting for this social feeling, Comte writes:

\begin{quotex}
There are three successive states of morality answering to the three principal stages of human life; the personal, the domestic, and the social stage. 

\end{quotex}
Between self-love and social feeling, Comte places domestic attachments on which ``the solution of the great moral problem depends." Comte elaborates:

\begin{quotex}
The love of his family leads Man out of his original state of Self-love and enables him to attain finally a sufficient measure of Social love. Every attempt on the part of the moral educator to call this last into immediate action, regardless of the intermediate stage, is to be condemned as utterly chimerical and profoundly injurious. Such attempts are regarded in the present day with far too favourable an eye. Far from being a sign of social progress, this would, if successful, be an immense step backwards … 

\end{quotex}
On the one hand, then, we see today — even more so than in Comte's — the attempt to induce an indiscriminate universal love in disregard for natural and organic familial relationships. These latter, instead, are under attack and ``family" becomes no more than a definition, fluidly and tendentiously redefined. On the other hand, the ``defenders" of family values limit themselves to an idealized image of the nuclear family and are unable to relate it to the larger social whole and probably would even reject the conclusion were they able to even see it clearly. Thus, true social feeling is forcibly separated from its natural and organic roots; in its place, there is only a legal fiction which can only be enforced by increasingly onerous laws and other social controls.

Comte identifies four familial relations that aid the transition from the individual to social feeling:

\begin{itemize}
\item Filial love 
\item Fraternal love 
\item Conjugal love 
\item Paternal love 
\end{itemize}
\paragraph{Filial Love}
Filial love is the affection of the child toward its parents. It is the starting point of moral education and the instinct of Continuity springs from it. The social consequence, then, is reverence for our ancestors which binds us to the whole past history of our people.

\paragraph{Fraternal Love}
Fraternal, or brotherly, love is the affection of child toward its brothers and sisters. From it, the instinct of Solidarity arises. Hence, we learn unity with contemporaries. Relationships become voluntary rather than involuntary.

\paragraph{Conjugal Love}
The second stage of moral education begins with conjugal love—``the highest type of all sympathetic instincts" — in which the ``perfect fullness of devotion is secured by the reciprocity and indissolubility of the bond."

\paragraph{Paternal Love}
From the conjugal relationship, paternal love comes next.

\begin{quotex}
It completes the training by which Nature prepares us for universal sympathy: for it teaches care for our successors; and thus it binds us to the Future, as filial love had bound us to the Past. 

\end{quotex}


\flrightit{Posted on 2008-11-30 by Cologero }
